
Note that when $h\rightarrow\infty$, the optimization objective in Equation \eqref{eq:local:reward:opt} is reduce to
\begin{equation} \label{eq:local:reward:opt:infty}
\begin{split}
    \max_{\pi}
    \mathop{\mathbb{E}}_{ 
    \substack{
    a_t \sim \pi(\cdot \mid s_t) \\ 
    s_{t+1} \sim M(\cdot \mid s_t, a_t)
    }}
    [\sum_{t} \gamma^t r_t]
\end{split}
\end{equation}
where $M = E(\alpha_k + \delta)$,
which is essentially deterministic in Equation \eqref{eq:det:opt} with evolution progression size $l_{k+1}$ equals to the neighborhood range $\delta$.


\begin{equation}
\begin{split}
    f(\alpha^\prime) & = (1+h\alpha^\prime)V^{\pi^*_{E(\alpha^\prime)},E(\alpha^\prime)}- \frac{1}{\delta} \int_{\beta=\alpha}^{\alpha+\delta} (1 + h \beta)
    V^{\pi^*_{E(\alpha^\prime)},E(\beta)}  \dd{\beta} \\
    & = \int_{\beta=\alpha}^{\alpha+\delta} (V^{\pi^*_{E(\alpha^\prime)},E(\alpha^\prime)} - V^{\pi^*_{E(\alpha^\prime)},E(\beta)} ) + h 
\end{split}
\end{equation}
is maximized at $\alpha^\prime$
\begin{equation}
\begin{split}
    & \lim_{\delta \rightarrow 0}
    \left[
    \mathop{\arg\max}_{\pi}
    \mathop{\mathbb{E}}_{ 
    \substack{
    \beta \sim U([\alpha, \alpha+\delta])
    }}
    \mathop{\mathbb{E}}_{ 
    \substack{(s_t, a_t) \sim \pi, E(\beta)
    }}
    \sum_{t} \gamma^t r_t (1 + h \beta)
    \right] \\
    & = \lim_{\delta \rightarrow 0}
    \left[
    \mathop{\arg\max}_{\pi}
    \mathop{\mathbb{E}}_{ 
    \substack{
    \beta \sim U([\alpha, \alpha+\delta])
    }}
    (1 + h \beta)
    \mathop{\mathbb{E}}_{ 
    \substack{(s_t, a_t) \sim \pi, E(\beta)
    }}
    \sum_{t} \gamma^t r_t 
    \right] \\
    & = \lim_{\delta \rightarrow 0}
    \left[
    \mathop{\arg\max}_{\pi}
    \mathop{\mathbb{E}}_{ 
    \substack{
    \beta \sim U([\alpha, \alpha+\delta])
    }}
    (1 + h \beta) V^{\pi, E(\beta)} 
    \right] \\
    & = \lim_{\delta \rightarrow 0}
    \left[
    \mathop{\arg\max}_{\pi}
    \int_{\beta=\alpha}^{\alpha+\delta} (1+h\beta)
    \right] \\
\end{split}
\end{equation}
    
    
    
    By the assumption in Equation \ref{eq:alpha:lipschitz:assumption},
\begin{equation}
\begin{split}
  \exists \delta > 0, \text{s.t. } \forall |\alpha^\prime - \alpha | < \delta, 
 ||E(\alpha)(s) - E(\alpha^\prime)(s)|| \le L_1 |\alpha - \alpha^\prime|,
 \forall s\in\mathcal{S} 
\end{split}
\end{equation}